\begin{table}[htb]
  \caption[Einfache lineare Regression von Religiosität auf psychisches Wohlbefinden]{\textit {Einfache lineare Regression von Religiosität auf psychisches Wohlbefinden}} 
  \label{H1_Regression}
  \centering
  \begin{adjustbox}{width=13cm} %{maxwidth=\textwidth}
  \small
  \begin{threeparttable}
      \begin{tabular}{lrrrrrr}
          \toprule
          & $B$ & $SE(B)$ & \textbeta & $t$ & \multicolumn{2}{c}{95\% Konfidenzintervall} \\
          \cmidrule(r){6-7}
          & & & & & Untere Grenze & Obere Grenze \\
          \midrule
          Konstante & 3.23 & 0.04 & & 75.03 & 3.14 & 3.31 \\
          Religiosität & 0.21 & 0.04 & .23* & 4.81 & 0.12 & 0.29 \\
          \bottomrule
      \end{tabular}
      \begin{tablenotes}
          \small
          \item \textit{Anmerkungen.} \(N\)~=~429; Wertebereich der Variable Religiosität -1.69 (\textit{gar nicht religiös}) bis 1.16 (\textit{sehr religiös}). \\ *$p<$~.001
      \end{tablenotes}
  \end{threeparttable}
  \end{adjustbox}
\end{table}

% psychisches Wohlbefinden 1 (\textit{die ganze Zeit}) bis 6 (\textit{zu keinem Zeitpunkt});